\chapter[Les tableaux]{
Les tableaux}
{
Dans ce chapitre nous étudions les tableaux, une structure qui peut
contenir plusieurs exemplaires de données similaires.}

\begin{center}
 [Warning: Image ignored] % Unhandled or unsupported graphics:
%\includegraphics[width=1.27cm,height=1.27cm]{log1-img/log1-img99}

\end{center}
\section{Utilité des tableaux}
{
Nous allons introduire la notion de tableau à partir d’un exemple dans
lequel l’utilité de cette structure de données apparaitra de façon
naturelle. }

{\sffamily\bfseries\scshape
Exemple : Statistiques de vente}

{\itshape
Un gérant d’une entreprise commerciale souhaite connaitre l’impact d’une
journée de promotion publicitaire sur la vente de dix de ses produits.
Pour ce faire, les numéros de ces produits (numérotés de 1 à 10 pour
simplifier) ainsi que les quantités vendues pendant cette journée de
promotion sont encodés au fur et à mesure de leurs ventes. En fin de
journée, le vendeur entrera la valeur 0 pour signaler la fin de
l’introduction des données. Ensuite, les statistiques des ventes seront
affichées.}

{
La démarche générale se décompose en trois parties :}

\liststyleListv
\begin{itemize}
\item {
le traitement de début de journée, qui consiste essentiellement à mettre
les compteurs des quantités vendues pour chaque produit à 0}
\item {
le traitement itératif durant toute la journée : au fur et à mesure des
ventes, il convient de les enregistrer, c’est-à-dire d’ajouter au
compteur des ventes d’un produit la quantité vendue de ce produit ; ce
traitement itératif s’interrompra lorsque la valeur 0 sera introduite}
\item {
le traitement final, consistant à communiquer les valeurs des compteurs
pour chaque produit.}
\end{itemize}
{
Vous trouverez sur la page suivante, une version possible de cet
algorithme.}

{
Il est clair, à la lecture de cet algorithme, qu’une simplification
d’écriture s’impose ! Et que ce passerait-il si le nombre de produits à
traiter était de 20 ou 100 ? Le but de l’informatique étant de
dispenser l’humain des tâches répétitives, le programmeur peut en
espérer autant de la part d’un langage de programmation ! La solution
est apportée par un nouveau type de variables : les \textbf{variables
indicées} ou \textbf{tableaux}.}

{\sffamily
\textbf{module} statistiquesVentesSansTableau()}


\bigskip

{\sffamily
\ \ cpt1, cpt2, cpt3, cpt4, cpt5 : entiers}

{\sffamily
\ \ cpt6, cpt7, cpt8, cpt9, cpt10 : entiers}

{\sffamily
\ \ numéroProduit, quantité : entiers}


\bigskip

{\sffamily
\ \ cpt1 ← 0}

{\sffamily
\ \ cpt2 ← 0}

{\sffamily
\ \ cpt3 ← 0}

{\sffamily
\ \ cpt4 ← 0}

{\sffamily
\ \ cpt5 ← 0}

{\sffamily
\ \ cpt6 ← 0}

{\sffamily
\ \ cpt7 ← 0}

{\sffamily
\ \ cpt8 ← 0}

{\sffamily
\ \ cpt9 ← 0}

{\sffamily
\ \ cpt10 ← 0}


\bigskip

{\sffamily
\textbf{\ \ écrire} {\textquotedbl}Introduisez le numéro du produit :
{\textquotedbl}}

{\sffamily
\textbf{\ \ lire} numéroProduit}


\bigskip

{\sffamily
\textbf{\ \ tant que} numéroProduit {\textgreater} 0 \textbf{faire}}


\bigskip

{\sffamily
\textbf{\ \ \ \ écrire} {\textquotedbl}Introduisez la quantité vendue :
{\textquotedbl}}

{\sffamily
\textbf{\ \ \ \ lire} quantité}


\bigskip

{\sffamily
\textbf{\ \ \ \ selon que} numéroProduit \textbf{vaut}}

{\sffamily
\ \ \ \ \ \ 1: cpt1 ← cpt1 + quantité}

{\sffamily
\ \ \ \ \ \ 2: cpt2 ← cpt2 + quantité}

{\sffamily
\ \ \ \ \ \ 3: cpt3 ← cpt3 + quantité}

{\sffamily
\ \ \ \ \ \ 4: cpt4 ← cpt4 + quantité}

{\sffamily
\ \ \ \ \ \ 5: cpt5 ← cpt5 + quantité}

{\sffamily
\ \ \ \ \ \ 6: cpt6 ← cpt6 + quantité}

{\sffamily
\ \ \ \ \ \ 7: cpt7 ← cpt7 + quantité}

{\sffamily
\ \ \ \ \ \ 8: cpt8 ← cpt8 + quantité}

{\sffamily
\ \ \ \ \ \ 9: cpt9 ← cpt9 + quantité}

{\sffamily
\ \ \ \ \ \ 10: cpt10 ← cpt10 + quantité}

{\sffamily
\ \ \ \ \textbf{fin selon}}


\bigskip

{\sffamily
\textbf{\ \ \ \ écrire} {\textquotedbl}Introduisez le numéro du produit
: {\textquotedbl}}

{\sffamily
\textbf{\ \ \ \ lire} numéroProduit}


\bigskip

{\sffamily\bfseries
\ \ fin tant}


\bigskip

{\sffamily
\textbf{\ \ }\textbf{écrire }{\textquotedbl}quantité vendue de produit 1
: {\textquotedbl}, cpt1}

{\sffamily
\textbf{\ \ écrire} {\textquotedbl}quantité vendue de produit 2 :
{\textquotedbl}, cpt2}

{\sffamily
\textbf{\ \ écrire} {\textquotedbl}quantité vendue de produit 3 :
{\textquotedbl}, cpt3}

{\sffamily
\textbf{\ \ écrire} {\textquotedbl}quantité vendue de produit 4 :
{\textquotedbl}, cpt4}

{\sffamily
\textbf{\ \ écrire} {\textquotedbl}quantité vendue de produit 5 :
{\textquotedbl}, cpt5}

{\sffamily
\textbf{\ \ écrire} {\textquotedbl}quantité vendue de produit 6 :
{\textquotedbl}, cpt6}

{\sffamily
\textbf{\ \ écrire} {\textquotedbl}quantité vendue de produit 7 :
{\textquotedbl}, cpt7}

{\sffamily
\textbf{\ \ écrire} {\textquotedbl}quantité vendue de produit 8 :
{\textquotedbl}, cpt8}

{\sffamily
\textbf{\ \ écrire} {\textquotedbl}quantité vendue de produit 9 :
{\textquotedbl}, cpt9}

{\sffamily
\textbf{\ \ écrire} {\textquotedbl}quantité vendue de produit 10 :
{\textquotedbl}, cpt10}


\bigskip

{\sffamily
\textbf{fin} \textbf{module}}


\bigskip


\bigskip


\bigskip

{
Au lieu d’avoir à manier dix compteurs distincts
(\textstyleCodeInsr{cpt1}, \textstyleCodeInsr{cpt2}, etc.), nous allons
envisager une seule «~grande~» variable cpt compartimentée en dix
«~sous-variables~» qui se distingueront les unes des autres par un
indice : \textstyleCodeInsr{cpt1} deviendrait ainsi
\textstyleCodeInsr{cpt[1]}, \textstyleCodeInsr{cpt2} deviendrait
\textstyleCodeInsr{cpt[2]}, et ainsi de suite jusqu’à
\textstyleCodeInsr{cpt10} qui deviendrait \textstyleCodeInsr{cpt[10]}.
}

\begin{center}
\tablehead{}
\begin{supertabular}{m{0.588cm}m{1.074cm}m{1.074cm}m{1.074cm}m{1.074cm}m{1.074cm}m{1.074cm}m{1.074cm}m{1.074cm}m{1.074cm}m{1.2989999cm}}
~
 &
\centering  \textstyleCodeInsr{cpt[1]} &
\centering  \textstyleCodeInsr{cpt[2]} &
\centering  \textstyleCodeInsr{cpt[3]} &
\centering  \textstyleCodeInsr{cpt[4]} &
\centering  \textstyleCodeInsr{cpt[5]} &
\centering  \textstyleCodeInsr{cpt[6]} &
\centering  \textstyleCodeInsr{cpt[7]} &
\centering  \textstyleCodeInsr{cpt[8]} &
\centering  \textstyleCodeInsr{cpt[9]} &
\centering\arraybslash 
\textstyleCodeInsr{cpt[10]}\\\hhline{~----------}
\multicolumn{1}{m{0.588cm}|}{\centering 
\textstyleCodeInsr{cpt}} &
\multicolumn{1}{m{1.074cm}|}{~
} &
\multicolumn{1}{m{1.074cm}|}{~
} &
\multicolumn{1}{m{1.074cm}|}{~
} &
\multicolumn{1}{m{1.074cm}|}{~
} &
\multicolumn{1}{m{1.074cm}|}{~
} &
\multicolumn{1}{m{1.074cm}|}{~
} &
\multicolumn{1}{m{1.074cm}|}{~
} &
\multicolumn{1}{m{1.074cm}|}{~
} &
\multicolumn{1}{m{1.074cm}|}{~
} &
\multicolumn{1}{m{1.2989999cm}|}{~
}\\\hhline{~----------}
\end{supertabular}
\end{center}
{
Un des intérêts de cette notation est la possibilité de faire apparaitre
une variable entre les crochets, par exemple
\textstyleCodeInsr{cpt[i]}, ce qui permet une grande économie de lignes
de code.}

{
Voici la version avec tableau.}

{\sffamily
\textbf{module} statistiquesVentesAvecTableau()}


\bigskip

{\sffamily
\ \ cpt : \textbf{tableau} [1 à 10] d’entiers}

{\sffamily
\ \ numéroProduit, quantité, i : entiers}


\bigskip

{\sffamily
\textbf{\ \ pour} i \textbf{de} 1 \textbf{à} 10 \textbf{faire}}

{\sffamily
\ \ \ \ cpt[i]\texttt{ }←\texttt{ }0}

{\sffamily\bfseries
\ \ fin pour}


\bigskip

{\sffamily
\textbf{\ \ écrire} {\textquotedbl}Introduisez le numéro du produit :
{\textquotedbl}}

{\sffamily
\textbf{\ \ lire} numéroProduit}


\bigskip

{\sffamily
\textbf{\ \ tant que} numéroProduit {\textgreater} 0 \textbf{faire}}


\bigskip

{\sffamily
\textbf{\ \ \ \ écrire} {\textquotedbl}Introduisez la quantité vendue :
{\textquotedbl}}

{\sffamily
\textbf{\ \ \ \ lire} quantité}


\bigskip

{\sffamily
\ \ \ \ cpt[numéroProduit\texttt{] }←\texttt{ }cpt[numéroProduit] +
quantité}


\bigskip

{\sffamily
\textbf{\ \ \ \ écrire} {\textquotedbl}Introduisez le numéro du produit
: {\textquotedbl}}

{\sffamily
\textbf{\ \ \ \ lire} numéroProduit}

{\sffamily\bfseries
\ \ fin tant}


\bigskip

{\sffamily
\textbf{\ \ pour} i \textbf{de} 1 \textbf{à} 10 \textbf{faire}}

{\sffamily
\textbf{\ \ \ \ écrire} {\textquotedbl}quantité vendue de
produit~{\textquotedbl}, i~, {\textquotedbl} : {\textquotedbl}, cpt[i]}

{\sffamily\bfseries
\ \ fin pour}


\bigskip

{\sffamily\bfseries
fin module}


\bigskip

\section{Définitions}
{
\textstyleWWPolicepardfaut{Un
}\textstyleWWPolicepardfaut{\textbf{tableau}}\textstyleWWPolicepardfaut{
est une suite d’éléments de même type portant tous le même nom mais se
distinguant les uns des autres par un indice. }}

\begin{center}
 [Warning: Image ignored] % Unhandled or unsupported graphics:
%\includegraphics[width=1.129cm,height=1.282cm]{log1-img/log1-img100}

\end{center}
{
\textstyleWWPolicepardfaut{L’}\textstyleWWPolicepardfaut{\textbf{indice}}\textstyleWWPolicepardfaut{
est un entier donnant la position d’un élément dans la suite. Cet
indice varie entre la position du premier élément et la position du
dernier élément, ces positions correspondant aux bornes de l’indice. }}

{
\textstyleWWPolicepardfaut{Notons qu'il
n'y a pas de «~trou~» : tous les éléments existent
entre le premier et le dernier indice.}}

{
\textstyleWWPolicepardfaut{La
}\textstyleWWPolicepardfaut{\textbf{taille}}\textstyleWWPolicepardfaut{
d’un tableau est le nombre (strictement positif ) de ses éléments.
Attention ! la taille d’un tableau ne peut pas être modifiée pendant
son utilisation.}}

\begin{center}
 [Warning: Image ignored] % Unhandled or unsupported graphics:
%\includegraphics[width=1.129cm,height=1.282cm]{log1-img/log1-img101}

\end{center}
{
\textstyleWWPolicepardfaut{Souvent on utilise un tableau plus grand que
le nombre utile de ses éléments. Seule une partie du tableau est
utilisée. On parle alors de taille
}\textstyleWWPolicepardfaut{\textbf{physique}}\textstyleWWPolicepardfaut{
(la taille maximale du tableau) et de taille
}\textstyleWWPolicepardfaut{\textbf{logique}}\textstyleWWPolicepardfaut{
(le nombre d'éléments effectivement utilisés)}}

\section{Notations}
{
Pour déclarer un tableau, on écrit :}

{\sffamily
nomTableau : \textstyleMotCl{tableau} [borneMin à borneMax] de
TypeElément}

\begin{center}
 [Warning: Image ignored] % Unhandled or unsupported graphics:
%\includegraphics[width=1.129cm,height=1.282cm]{log1-img/log1-img102}

\end{center}
{
où \textstyleCodeInsr{TypeElément} est le type des éléments que l’on
trouvera dans le tableau. Les éléments sont d’un des types élémentaires
vus précédemment (entier, réel, booléen, chaine, caractère) ou encore
des variables structurées. À ce propos, remarquons aussi
qu'un tableau peut être un champ
d'une structure. D'autres
possibilités apparaitront lors de l'étude de
l'orienté objet. }

{
Les bornes apparaissant dans la déclaration sont des constantes ou des
paramètres ayant une valeur connue lors de la déclaration. Une fois un
tableau déclaré, seuls les éléments d’indice compris entre
\textstyleCodeInsr{borneMin} et \textstyleCodeInsr{borneMax} peuvent
être utilisés. Par exemple, si on déclare :}

{\sffamily
tabEntiers : \textstyleMotCl{tableau} [1 à 100] d’entiers}

{
Il est interdit d’utiliser \textstyleCodeInsr{tabEntiers[0]} ou
\textstyleCodeInsr{tabEntiers[101]}. De plus, chaque élément
\textstyleCodeInsr{tabEntiers[i]} (avec 1 ${\leq}$ i ${\leq}$ 100) doit
être manié avec la même précaution qu’une variable simple, c’est-à-dire
qu’on ne peut utiliser un élément du tableau qui n’aurait pas été
préalablement affecté ou initialisé. }

{
N.B. : Il n'est pas interdit de prendre 0 pour la borne
inférieure ou même d'utiliser des bornes négatives
(par exemple : \textstyleCodeInsr{tabTempératures[ -20 à 50 ]}).}

{
En Java, un tableau est défini par sa taille n et les bornes sont
automatiquement 0 et n-1. Ce n'est pas le cas en
Logique où on a plus de liberté dans le choix des bornes.}

\begin{center}
 [Warning: Image ignored] % Unhandled or unsupported graphics:
%\includegraphics[width=1.323cm,height=1.323cm]{log1-img/log1-img103}

\end{center}
\section{Tableau statique vs tableau dynamique}
{
\textstyleWWPolicepardfaut{Les tableaux étudiés
jusqu'ici sont dit
}\textstyleWWPolicepardfaut{\textbf{statiques}}\textstyleWWPolicepardfaut{.
Un tableau est dit
}\textstyleWWPolicepardfaut{\textbf{statique}}\textstyleWWPolicepardfaut{
si sa taille est connue lors de l’écriture du programme}}

{
\textstyleWWPolicepardfaut{Un tableau est dit
}\textstyleWWPolicepardfaut{\textbf{dynamique}}\textstyleWWPolicepardfaut{
si sa taille n’est connue qu’à l’exécution du programme (elle est
calculée, donnée par l’utilisateur, lue dans un fichier de
configuration, fournie par une autre partie du programme, …)}}

{
\textstyleWWPolicepardfaut{\textbf{Exemples}}}

\liststyleListv
\begin{itemize}
\item {
\textstyleWWPolicepardfaut{Reprenons l’exemple ci-dessus. S’il y a
toujours exactement 10 produits, alors }\textstyleWWPolicepardfaut{on
peut utiliser un tableau statique. Dans la pratique, il est probable
que ce }\textstyleWWPolicepardfaut{nombre puisse évoluer au cours de
l’histoire de l’entreprise.}}
\item {
\textstyleWWPolicepardfaut{Si on désire stocker le chiffre d’affaire
d’une entreprise pour une année donnée, mois par mois, un tableau de
taille 12 est suffisant (tableau statique)}}
\end{itemize}
{
\textstyleWWPolicepardfaut{Certains langages (c’est le cas de Cobol) ne
permettent que des tableaux statiques. Dans ce cas, il faudra souvent
imposer une }\textstyleWWPolicepardfaut{\textbf{taille
maximale}}\textstyleWWPolicepardfaut{ au tableau. Il sera également
nécessaire d’utiliser une variable pour retenir la taille utile (ou
logique) du tableau, à savoir la partie du tableau réellement utilisée.
Bien sûr cette solution entraine souvent une perte de place mémoire due
à une réservation inutilement grande (ou, pire, une saturation du
tableau qui n’aura pas été prévu assez grand)}}

{
\textstyleWWPolicepardfaut{\textbf{Exemple}}}

{
\textstyleWWPolicepardfaut{Reprenons encore l’exemple de la vente de
produits. Si on ne dispose pas de tableau dynamique, on peut réserver
un tableau de taille 1000 par exemple (si on est sûr qu’on ne vendra
pas plus de 1000 produits différents) et ajouter une variable entière
(}\textstyleCodeInsr{nbProduits}\textstyleWWPolicepardfaut{) pour
retenir le nombre exact de produits différents} actuellement en vente.}

{
Pour déclarer un tableau dynamique, on sépare la déclaration du tableau
(où la taille n’est pas donnée) de sa création effective~(où on donne
la taille)}

{\sffamily
nomTableau : \textstyleMotCl{tableau} de TypeElément}

\begin{center}
 [Warning: Image ignored] % Unhandled or unsupported graphics:
%\includegraphics[width=1.129cm,height=1.282cm]{log1-img/log1-img104}

\end{center}
{\sffamily
… }

{\sffamily
// le code peut déterminer ici les bornes}

{\sffamily
nomTableau~ ←\texttt{ }\textstyleMotCl{nouveau} \textstyleMotCl{tableau}
[ borneMin à borneMax ] de TypeElément}

{
où \textstyleCodeInsr{borneMin} et \textstyleCodeInsr{borneMax} sont des
expressions entières quelconques.}

\section{Tableau et paramètres}
{
Un tableau peut être passé en paramètre à un module mais qu’en est-il de
sa taille ? Il serait utile de pouvoir appeler le même module avec des
tableaux de tailles différentes. Pour permettre cela, la taille du
tableau reçu en paramètre est déclarée avec une variable (qui peut être
considéré comme un paramètre entrant). }

{
Exemple~:}

{\sffamily
\textstyleMotCl{module} brol( tab : \textstyleMotCl{tableau} [1 à n]
d'entiers )}

{\sffamily
\ \ \textstyleMotCl{écrire} {\textquotedbl}J’ai reçu un tableau de
{\textquotedbl}, n, {\textquotedbl}~éléments.{\textquotedbl}}

{\sffamily
\textstyleMotCl{fin} \textstyleMotCl{module}}

{
Ce \textstyleCodeInsr{n} va prendre la taille précise du tableau
(statique ou dynamique) utilisé à chaque appel et peut être utilisé
dans le corps du module. Bien sûr il s’agit là de la taille physique du
tableau. Si une partie seulement du tableau doit être traitée, il
convient de passer également la taille logique en paramètre.}

{
Notons qu’on admet également qu’un module retourne un tableau. Par
contre, on ne peut pas lire ou écrire un tableau en une seule
instruction ; il faut des instructions de lecture ou
d'écriture individuelles pour chacun de ses éléments.}

\section{Exercices}
{\sffamily\bfseries
Exercices de base}

\liststyleExercice
\begin{enumerate}
\item {\sffamily\bfseries
Somme}
\end{enumerate}
{
Écrire un module qui reçoit en paramètre le tableau
\textstyleCodeInsr{tabEnt} de \textstyleCodeInsr{n} entiers et qui
retourne la somme de ses éléments.}

\begin{center}
 [Warning: Image ignored] % Unhandled or unsupported graphics:
%\includegraphics[width=1.134cm,height=1.134cm]{log1-img/log1-img105}

\end{center}
\liststyleExercice
\setcounter{saveenum}{\value{enumi}}
\begin{enumerate}
\setcounter{enumi}{\value{saveenum}}
\item {\sffamily\bfseries
Maximum/minimum}
\end{enumerate}
{
Écrire un module qui reçoit en paramètre le tableau
\textstyleCodeInsr{tabEnt} de \textstyleCodeInsr{n} entiers et qui
retourne la plus grande valeur de ce tableau. Idem pour le minimum.}

\begin{center}
 [Warning: Image ignored] % Unhandled or unsupported graphics:
%\includegraphics[width=1.134cm,height=1.134cm]{log1-img/log1-img106}

\end{center}
\liststyleExercice
\setcounter{saveenum}{\value{enumi}}
\begin{enumerate}
\setcounter{enumi}{\value{saveenum}}
\item {\sffamily\bfseries
Indice du maximum/minimum}
\end{enumerate}
{
Écrire un module qui reçoit en paramètre le tableau
\textstyleCodeInsr{tabEnt} de \textstyleCodeInsr{n} entiers et qui
retourne l’indice de l’élément contenant la plus grande valeur de ce
tableau. 
En cas d’ex-æquo, c’est l’indice le plus petit qui sera renvoyé.

Que faut-il changer pour renvoyer l’indice le plus grand ?
Et pour retourner l’indice du minimum ? 
réécrire l’algorithme de l’exercice précédent en utilisant celui-ci.}

\begin{center}
 [Warning: Image ignored] % Unhandled or unsupported graphics:
%\includegraphics[width=1.134cm,height=1.134cm]{log1-img/log1-img107}

\end{center}
\liststyleExercice
\setcounter{saveenum}{\value{enumi}}
\begin{enumerate}
\setcounter{enumi}{\value{saveenum}}
\item {\sffamily\bfseries
Nombre d'éléments d'un tableau}
\end{enumerate}
{
Écrire un module qui reçoit en paramètre le tableau
\textstyleCodeInsr{tabEnt} de \textstyleCodeInsr{n} entiers et qui
retourne le nombre d’éléments du tableau. }

\liststyleExercice
\setcounter{saveenum}{\value{enumi}}
\begin{enumerate}
\setcounter{enumi}{\value{saveenum}}
\item {\sffamily\bfseries
La présence d'un 0}
\end{enumerate}
{
Écrire un module qui reçoit en paramètre le tableau
\textstyleCodeInsr{tabEnt} de \textstyleCodeInsr{n} entiers et qui
retourne un booléen indiquant s'il contient au moins
un 0. }

\liststyleExercice
\setcounter{saveenum}{\value{enumi}}
\begin{enumerate}
\setcounter{enumi}{\value{saveenum}}
\item {\sffamily\bfseries
Plus grand écart absolu}
\end{enumerate}
{
Écrire un module qui reçoit en paramètre le tableau
\textstyleCodeInsr{tabEnt} de \textstyleCodeInsr{n} entiers et qui
retourne le plus grand écart absolu entre deux éléments consécutifs de
ce tableau.
Et si on veut le plus petit écart ?}

\liststyleExercice
\setcounter{saveenum}{\value{enumi}}
\begin{enumerate}
\setcounter{enumi}{\value{saveenum}}
\item {\sffamily\bfseries
Remplacement de valeurs}
\end{enumerate}
{
Écrire un module qui reçoit en paramètre le tableau
\textstyleCodeInsr{tabEnt} de \textstyleCodeInsr{n} entiers et qui
remplace dans ce tableau toutes les valeurs multiples de 3 par 0.}

\liststyleExercice
\setcounter{saveenum}{\value{enumi}}
\begin{enumerate}
\setcounter{enumi}{\value{saveenum}}
\item {\sffamily\bfseries
Tableau ordonné ?}
\end{enumerate}
{
Écrire un module qui reçoit en paramètre le tableau
\textstyleCodeInsr{tabEnt} de \textstyleCodeInsr{n} entiers et qui
vérifie si ce tableau est ordonné (strictement) croissant sur les
valeurs. Le module retournera \textbf{vrai} si le tableau est ordonné,
\textbf{faux} sinon. }

\begin{center}
 [Warning: Image ignored] % Unhandled or unsupported graphics:
%\includegraphics[width=1.134cm,height=1.134cm]{log1-img/log1-img108}

\end{center}
\liststyleExercice
\setcounter{saveenum}{\value{enumi}}
\begin{enumerate}
\setcounter{enumi}{\value{saveenum}}
\item {\sffamily\bfseries
Position des minimums}
\end{enumerate}
{
Écrire un module qui reçoit en paramètre le tableau
\textstyleCodeInsr{tabEnt} de \textstyleCodeInsr{n} entiers et qui
affiche le ou les indice(s) des éléments contenant la valeur minimale
du tableau.}

\liststyleNumberingv
\begin{enumerate}
\item {
écrire une première version «~classique~» avec deux parcours de tableau}
\item {
écrire une deuxième version qui ne parcourt qu’une seule fois tabEnt en
stockant dans un deuxième tableau les indices des plus petits éléments
rencontrés (ce tableau étant à chaque fois réinitialisé lorsqu’un
nouveau minimum est rencontré)}
\begin{center}
 [Warning: Image ignored] % Unhandled or unsupported graphics:
%\includegraphics[width=1.134cm,height=1.134cm]{log1-img/log1-img109}

\end{center}
\end{enumerate}
\liststyleExercice
\begin{enumerate}
\item {\sffamily\bfseries
Renverser un tableau}
\end{enumerate}
{
Écrire un module qui reçoit en paramètre le tableau
\textstyleCodeInsr{tabEnt} de \textstyleCodeInsr{n} entiers, et qui
«~renverse~» ce tableau, c’est-à-dire qui permute le premier élément
avec le dernier, le deuxième élément avec l’avant-dernier et ainsi de
suite.}

\begin{center}
 [Warning: Image ignored] % Unhandled or unsupported graphics:
%\includegraphics[width=1.134cm,height=1.134cm]{log1-img/log1-img110}

\end{center}
\liststyleExercice
\setcounter{saveenum}{\value{enumi}}
\begin{enumerate}
\setcounter{enumi}{\value{saveenum}}
\item {\sffamily\bfseries
Tableau symétrique ?}
\end{enumerate}
{
Écrire un module qui reçoit en paramètre le tableau
\textstyleCodeInsr{tabEnt} de \textstyleCodeInsr{n} entiers et qui
vérifie si ce tableau est symétrique, c’est-à-dire si le premier
élément est identique au dernier, le deuxième à l’avant-dernier et
ainsi de suite.}

\liststyleExercice
\setcounter{saveenum}{\value{enumi}}
\begin{enumerate}
\setcounter{enumi}{\value{saveenum}}
\item {\sffamily\bfseries
Cumul des ventes}
\end{enumerate}
{
Soit le tableau \textstyleCodeInsr{ventes[1~à 12]} d’entiers où le
premier élément contient le montant total des ventes pour janvier, le
second pour février, et ainsi de suite jusqu'à
décembre. Écrire l’algorithme qui reçoit ce tableau en paramètre, et
renvoie le tableau \textstyleCodeInsr{cumul[1~à 12]} ; chaque élément
de ce tableau devra contenir le cumul de toutes les ventes depuis le
début de l’année jusqu’au mois correspondant à
l'indice de l’élément en question. Le dernier élément
de cumul devra donc contenir le total des ventes de l’année. }

\liststyleExercice
\setcounter{saveenum}{\value{enumi}}
\begin{enumerate}
\setcounter{enumi}{\value{saveenum}}
\item {\sffamily\bfseries
Occurrence des chiffres}
\end{enumerate}
{
Écrire un module qui reçoit un nombre entier positif ou nul en paramètre
et qui affiche pour chacun de ses chiffres le nombre de fois qu’il
apparait dans ce nombre. Ainsi, pour le nombre 10502851125, l’affichage
mentionnera que le chiffre 0 apparait 2 fois, 1 apparait 3 fois, 2
apparait 2 fois, 5 apparait 3 fois et 8 apparait une fois (l’affichage
ne mentionnera donc pas les chiffres qui n’apparaissent pas).}

\begin{center}
 [Warning: Image ignored] % Unhandled or unsupported graphics:
%\includegraphics[width=1.134cm,height=1.134cm]{log1-img/log1-img111}

\end{center}
\liststyleExercice
\setcounter{saveenum}{\value{enumi}}
\begin{enumerate}
\setcounter{enumi}{\value{saveenum}}
\item {\sffamily\bfseries
Palindrome}
\end{enumerate}
{
Soit le tableau \textstyleCodeInsr{phrase[1 à n]} dont chaque élément
est un caractère alphabétique, le caractère d’espacement ou un
caractère de ponctuation. Dans ce tableau, des lettres qui se suivent
constituent un mot. Ces mots sont séparés les uns des autres par un ou
plusieurs caractères d’espacement ou de ponctuation. Écrire un module
qui reçoit en paramètre ce tableau et qui vérifie si la phrase formée
par les mots de ce tableau est un palindrome. Le résultat sera
communiqué par le renvoi d’une valeur booléenne. Pour rappel, une
phrase palindrome est une phrase qui peut se lire dans les deux sens
sans tenir compte des espacements et de la ponctuation. }

{
Exemple du contenu du tableau phrase :}


\bigskip

\begin{center}
\tablehead{}
\begin{supertabular}{|m{0.476cm}|m{0.476cm}|m{0.476cm}|m{0.476cm}|m{0.476cm}|m{0.476cm}|m{0.476cm}|m{0.476cm}|m{0.476cm}|m{0.476cm}|m{0.476cm}|m{0.476cm}|m{0.476cm}|m{0.476cm}|m{0.476cm}|m{0.476cm}|m{0.476cm}|m{0.476cm}|m{0.476cm}|m{0.476cm}|m{0.5cm}|}
\hline
\centering  A &
~
 &
\centering  L &
\centering  ' &
\centering  E &
\centering  T &
\centering  A &
\centering  P &
\centering  E &
\centering  , &
~
 &
~
 &
\centering  E &
\centering  P &
\centering  A &
\centering  T &
\centering  E &
\centering  {}- &
\centering  L &
\centering  A &
\centering\arraybslash  !\\\hline
\end{supertabular}
\end{center}
{
\textbf{Aide} : il est permis d’utiliser les modules suivants :}


\bigskip

{\sffamily
\textbf{module} estEspace(car : caractère) → booléen }

{\sffamily
\ \ \ \ \ \ \ \ // retourne vrai si car est un caractère d’espacement}

{\sffamily
\textbf{module} estPonctuation(car : caractère) → booléen }

{\sffamily
\ \ \ \ \ \ \ \ // retourne vrai si car est un caractère de ponctuation.
}


\bigskip

\liststyleExercice
\setcounter{saveenum}{\value{enumi}}
\begin{enumerate}
\setcounter{enumi}{\value{saveenum}}
\item {\sffamily\bfseries
Moyenne d'éléments}
\end{enumerate}
{
Soit le tableau \textstyleCodeInsr{tabEnt} contenant
\textstyleCodeInsr{n} entiers \textbf{différents}. Écrire l’algorithme
qui calcule et affiche la moyenne des éléments situés entre les valeurs
minimale et maximale du tableau (ces deux valeurs participant au calcul
de cette moyenne). Dans l’exemple ci-dessous, le maximum est 85, le
minimum 5 et la moyenne des 8 éléments concernés est 21. Exploitez
l'exercice 3 pour réaliser cet algorithme.}

\begin{center}
 [Warning: Image ignored] % Unhandled or unsupported graphics:
%\includegraphics[width=1.134cm,height=1.134cm]{log1-img/log1-img112}

\end{center}
\begin{center}
\tablehead{}
\begin{supertabular}{|m{0.47000003cm}|m{0.47000003cm}|m{0.47000003cm}|m{0.47000003cm}|m{0.315cm}|m{0.31700003cm}|m{0.47000003cm}|m{0.47000003cm}|m{0.31700003cm}|m{0.47000003cm}|m{0.47000003cm}|m{0.47000003cm}|m{0.47000003cm}|m{0.47000003cm}|m{0.47000003cm}|m{0.47000003cm}|m{0.47000003cm}|m{0.47000003cm}|m{0.47000003cm}|m{0.48800004cm}|}
\hline
\centering  12 &
\centering  85 &
\centering  21 &
\centering  17 &
\centering  8 &
\centering  6 &
\centering  10 &
\centering  16 &
\centering  5 &
\centering  74 &
\centering  64 &
\centering  29 &
\centering  41 &
\centering  11 &
\centering  73 &
\centering  72 &
\centering  28 &
\centering  66 &
\centering  55 &
\centering\arraybslash  44\\\hline
\end{supertabular}
\end{center}

\bigskip

\liststyleExercice
\setcounter{saveenum}{\value{enumi}}
\begin{enumerate}
\setcounter{enumi}{\value{saveenum}}
\item {\sffamily\bfseries
OXO}
\end{enumerate}
{
Soit le tableau \textstyleCodeInsr{Oxo} contenant \textstyleCodeInsr{n}
caractères. Chaque élément est le caractère ‘O’ ou ‘X’. Écrire
l’algorithme qui permet de compter et d’afficher le nombre de séquences
distinctes des valeurs consécutives ‘O’, ‘X’, ‘O’. Lorsqu’un ‘O’ fait
partie de deux séquences, on ne comptabilise qu’une seule séquence
(ainsi le ‘O’ en position 6 ci-dessous est le dernier d’une séquence et
le premier d’une autre : on ne comptabilisera qu’une seule séquence
pour ces deux séquences qui se chevauchent). Dans l’exemple ci-dessous,
il y a donc 3 séquences à comptabiliser :}


\bigskip

\begin{center}
\tablehead{}
\begin{supertabular}{m{0.47000003cm}m{0.47000003cm}m{0.47000003cm}m{0.47000003cm}m{0.35400003cm}m{0.35400003cm}m{0.47000003cm}m{0.47000003cm}m{0.35400003cm}m{0.47000003cm}m{0.47000003cm}m{0.47000003cm}m{0.47000003cm}m{0.47000003cm}m{0.47000003cm}m{0.47000003cm}m{0.47000003cm}m{0.47000003cm}m{0.47000003cm}m{0.48800004cm}}
\centering \itshape 1 &
\centering \itshape 2 &
\centering \itshape 3 &
\centering \itshape 4 &
\centering \itshape 5 &
\centering \itshape 6 &
\centering \itshape 7 &
\centering \itshape 8 &
\centering \itshape 9 &
\centering \itshape 10 &
\centering \itshape 11 &
\centering \itshape 12 &
\centering \itshape 13 &
\centering \itshape 14 &
\centering \itshape 15 &
\centering \itshape 16 &
\centering \itshape 17 &
\centering \itshape 18 &
\centering \itshape 19 &
\centering\arraybslash \itshape 20\\\hline
\multicolumn{1}{|m{0.47000003cm}|}{\centering 
O} &
\multicolumn{1}{m{0.47000003cm}|}{\centering  X}
&
\multicolumn{1}{m{0.47000003cm}|}{\centering  X}
&
\multicolumn{1}{m{0.47000003cm}|}{\centering  O}
&
\multicolumn{1}{m{0.35400003cm}|}{\centering  X}
&
\multicolumn{1}{m{0.35400003cm}|}{\centering  O}
&
\multicolumn{1}{m{0.47000003cm}|}{\centering  X}
&
\multicolumn{1}{m{0.47000003cm}|}{\centering  O}
&
\multicolumn{1}{m{0.35400003cm}|}{\centering  X}
&
\multicolumn{1}{m{0.47000003cm}|}{\centering  O}
&
\multicolumn{1}{m{0.47000003cm}|}{\centering  O}
&
\multicolumn{1}{m{0.47000003cm}|}{\centering  O}
&
\multicolumn{1}{m{0.47000003cm}|}{\centering  O}
&
\multicolumn{1}{m{0.47000003cm}|}{\centering  X}
&
\multicolumn{1}{m{0.47000003cm}|}{\centering  X}
&
\multicolumn{1}{m{0.47000003cm}|}{\centering  O}
&
\multicolumn{1}{m{0.47000003cm}|}{\centering  X}
&
\multicolumn{1}{m{0.47000003cm}|}{\centering  O}
&
\multicolumn{1}{m{0.47000003cm}|}{\centering  O}
&
\multicolumn{1}{m{0.48800004cm}|}{\centering\arraybslash
 X}\\\hline
\end{supertabular}
\end{center}

\bigskip

\liststyleExercice
\setcounter{saveenum}{\value{enumi}}
\begin{enumerate}
\setcounter{enumi}{\value{saveenum}}
\item {\sffamily\bfseries
Les doublons}

{
Vérifiez si un tableau contient au moins 2 éléments égaux.}

\begin{center}
 [Warning: Image ignored] % Unhandled or unsupported graphics:
%\includegraphics[width=1.134cm,height=1.134cm]{log1-img/log1-img113}

\end{center}

\bigskip
\item {\sffamily\bfseries
Mastermind}
\end{enumerate}
{
Dans le jeu du Mastermind, un joueur A doit trouver une combinaison de k
pions de couleur, choisie et tenue secrète par un autre joueur B. Cette
combinaison peut contenir éventuellement des pions de même couleur. À
chaque proposition du joueur A, le joueur B indique le nombre de pions
de la proposition qui sont corrects et bien placés et le nombre de
pions corrects mais mal placés. }

\begin{center}
 [Warning: Image ignored] % Unhandled or unsupported graphics:
%\includegraphics[width=1.134cm,height=1.134cm]{log1-img/log1-img114}

\end{center}
{
Pour implémenter une simulation de ce jeu, on utilise le type Couleur,
dont les valeurs possibles sont les couleurs des pions utilisés.
(Attention, le nombre exact de couleurs n’est pas précisé.) Les seules
manipulations permises avec ce type sont la comparaison (tester si deux
couleurs sont identiques ou non) et l’affectation (affecter le contenu
d’une variable de type Couleur à une autre variable de ce type). Les
propositions du joueur A, ainsi que la combinaison secrète du joueur B
sont contenues dans des tableaux de k composantes de type Couleur.}

{
Écrire le module suivant qui renvoie dans les variables
\textstyleCodeInsr{bienPlacés} et \textstyleCodeInsr{malPlacés}
respectivement le nombre de pions bien placés et mal placés dans la
«~proposition~» du joueur A en la comparant à la «~solution~» cachée du
joueur B.}

{\sffamily
\textbf{module} testerProposition( \ \ proposition
{↓}, }

{\sffamily
\ \ \ \ \ \ \ \ \ \ \ \ solution {↓} :
\textbf{tableau} [1 à k ] de Couleur, }

{\sffamily
\ \ \ \ \ \ \ \ \ \ \ \ bienPlacés ↑, malPlacés ↑~: entiers)}


\bigskip


\bigskip

